\chapter{凸优化、KKT 与对偶}
\label{ch:convex-optimization}

优化器给出的“最优”只有在目标、约束和数值证据都清楚时才有意义。本章先用几何语言建立凸性，再把约束资格、KKT 条件和对偶价格连成一条证明链，最后用一个二次规划和两资产组合说明如何从数值解恢复可检查的证书。读者应始终区分“找到了一个驻点”“找到了可行点”和“证明了全局最优”。

这一章的手算例子故意足够小，可以在纸上完成；失败例子则展示不可行、无界、资格条件缺失和尺度不良时，为什么不能把求解器状态当作数学结论。

\section{凸几何把局部证据升级为全局证据}

集合 $C$ 为凸集，若任意 $x,y\in C$ 与 $\theta\in[0,1]$ 都有 $\theta x+(1-\theta)y\in C$。函数 $f:C\to\mathbb R$ 为凸函数，若
\[
f(\theta x+(1-\theta)y)\le \theta f(x)+(1-\theta)f(y).
\]
等价地，函数上图 $\operatorname{epi}f=\{(x,t):x\in C,t\ge f(x)\}$ 是凸集。可微时一阶下界
\[
f(y)\ge f(x)+\nabla f(x)^T(y-x)
\]
刻画凸性；二阶可微时，凸定义域上 $\nabla^2f(x)\succeq0$ 与凸性等价。若 $\nabla^2f(x)\succeq mI$、$m>0$，则 $f$ 为 $m$-强凸，若最优解存在则唯一且目标差能控制距离；严格凸保证至多一个最优解，却未必给出统一曲率常数。

不可微时，用次梯度集合
\[
\partial f(x)=\{g:f(y)\ge f(x)+g^T(y-x),\ \forall y\}
\]
替代单一梯度。例如 $\partial|x|_{x=0}=[-1,1]$，故 $0\in\partial|x|_{x=0}$ 直接证明零点最优。闭凸集上的 Euclidean 投影
\[
\Pi_C(z)=\arg\min_{x\in C}\tfrac12\|x-z\|_2^2
\]
唯一，并满足非扩张性 $\|\Pi_C(a)-\Pi_C(b)\|\le\|a-b\|$。投影保持可行不等于原目标已经最优。

\begin{proposition}[凸目标的一阶最优性]
设 $C$ 是凸集，$f$ 在 $C$ 上可微且凸。点 $x^*\in C$ 是全局最小点，当且仅当
\[
  \nabla f(x^*)^T(y-x^*)\ge0,\qquad\forall y\in C.
\]
\end{proposition}
\begin{proof}
若 $x^*$ 最优，沿着线段 $x^*+\theta(y-x^*)$ 的右导数不能为负，故得到不等式。反过来，由凸函数的一阶下界，
\[
f(y)\ge f(x^*)+\nabla f(x^*)^T(y-x^*)\ge f(x^*).
\]
因此一个局部的方向性条件在凸问题中已经是全局证书；非凸时同样的驻点只说明没有一阶下降方向，不能排除更远处的更小值。
\end{proof}

强凸性提供的不只是唯一性。若
\[
 f(y)\ge f(x)+\nabla f(x)^T(y-x)+\frac m2\|y-x\|_2^2,
 \qquad m>0,
\]
且 $x^*$ 满足一阶最优性，则
\[
  f(y)-f(x^*)\ge \frac m2\|y-x^*\|_2^2.
\]
所以目标值接近最优时，变量本身也必须接近；普通严格凸函数没有统一的 $m$，可能只能保证解唯一而不能给出这样的误差换算。

\MFQLead{标准约束形式与建模记录}

考虑
\[
\min_x f(x),\qquad g_i(x)\le0,\qquad h_j(x)=0.
\]
凸优化要求 $f,g_i$ 凸而 $h_j$ 仿射。可行域、变量单位和每条约束方向必须先声明；$g_i(x^*)=0$ 的约束称为活跃约束。若收益按日、协方差按年、成本按基点却直接相加，形式即使凸也没有一致经济单位。等式暴露约束、预算约束和盒约束若互相冲突，则问题不可行；求解器的默认容差可能只返回近似违反。

Lagrangian 为
\[
L(x,\lambda,\nu)=f(x)+\sum_i\lambda_i g_i(x)+\sum_j\nu_jh_j(x),
\qquad \lambda_i\ge0.
\]
KKT 是四本同时对平的账：原可行、对偶可行、驻点
\[
0\in\partial f(x^*)+\sum_i\lambda_i^*\partial g_i(x^*)
+\sum_j\nu_j^*\nabla h_j(x^*),
\]
以及互补松弛 $\lambda_i^*g_i(x^*)=0$。一般非凸问题中，KKT 只是满足约束资格的局部最优必要条件；凸问题中，任何 KKT 点都是全局最优。

约束资格不能省略。LICQ 要求活跃不等式与等式梯度线性无关；Slater 条件要求凸问题存在严格满足所有凸不等式并满足等式的点。Slater 常用来保证强对偶与乘子存在，但不是所有问题唯一可用的资格条件。

\begin{theorem}[凸问题的 KKT 充分性]
设 $f$ 和 $g_i$ 凸、$h_j$ 仿射。若 $(x^*,\lambda^*,\nu^*)$ 满足原可行、对偶可行、驻点和互补松弛，则 $x^*$ 是原问题的全局最优点；若原问题已有最优解，并满足适用的 Slater 条件和次梯度正则性，则强对偶成立且最优解可配上这样的乘子。
\end{theorem}
\begin{proof}
对任意可行 $x$，取驻点中的次梯度 $s_f,s_i$，有
\[
 f(x)-f(x^*)\ge s_f^T(x-x^*),\qquad
 g_i(x)-g_i(x^*)\ge s_i^T(x-x^*).
\]
把驻点关系代入第一式，并用 $h_j(x)=h_j(x^*)=0$，得到
\[
 f(x)-f(x^*)\ge
 -\sum_i\lambda_i^*s_i^T(x-x^*)
 \ge\sum_i\lambda_i^*\{g_i(x^*)-g_i(x)\}\ge0.
\]
最后一步使用 $\lambda_i^*\ge0$、$g_i(x)\le0$ 和互补松弛。Slater 的作用不是证明这个不等式，而是保证原问题和对偶问题之间没有间隙，并使乘子可作为约束边际价值解释。
\end{proof}

\section{对偶、KKT 与影子价格}

对偶函数定义为 $q(\lambda,\nu)=\inf_xL(x,\lambda,\nu)$，它作为仿射函数族的逐点下确界总是凹。若 $x$ 原可行、$\lambda\ge0$，则
\[
q(\lambda,\nu)\le L(x,\lambda,\nu)\le f(x),
\]
所以任意对偶可行值都是原最小化问题的下界，这就是弱对偶。最大化该下界得到对偶问题；原最优值 $p^*$ 与对偶最优值 $d^*$ 满足 $d^*\le p^*$。凸性加适当资格条件（例如 Slater）给 $p^*=d^*$，但“凸”本身不能自动排除不可行、无界或最优值不取到。

把约束写成 $g(x)\le b$ 时，最优乘子描述收紧右端项的局部价值；符号取决于约束写法。影子价格是最优值函数的导数或次梯度，不是资产收益预测，也只在活跃集和正则性不突变的局部范围内可靠。

\MFQLead{手算二次规划：从 KKT 到对偶证书}

考虑
\[
\min_{x,y}\ \tfrac12(x^2+y^2),\qquad x+y\ge1,\qquad x,y\ge0.
\]
把主约束写成 $1-x-y\le0$。解为正，所以非负约束不活跃、对应乘子为零。主约束若不活跃，驻点会给 $x=y=0$，却不可行；故它必须活跃。Lagrangian
\[
L=\tfrac12(x^2+y^2)+\lambda(1-x-y)
\]
的驻点给 $x=y=\lambda$，再由 $x+y=1$ 得手算最优解为 $(0.5,0.5)$，最优乘子 $\lambda^*=0.5$。

% Generated by tools/build_figures.py; edit figures/figure-specs.json instead.
\MFQFigure{tex/figures/generated/upper-ch13-kkt-geometry.pdf}{图示问题的最优点位于目标等高线与活跃约束的切点；KKT 乘子描述目标梯度与约束法向量的关系。}{upper-ch13-kkt-geometry}


固定 $\lambda\ge0$ 对 $x,y$ 极小化，代入 $x=y=\lambda$，得到
\[
q(\lambda)=\lambda-\lambda^2.
\]
其最大点为 $\lambda=0.5$。原问题与对偶目标均为 $0.25$，对偶间隙为零；点 $(1,1)$ 是 Slater 点，强对偶并非偶然数值巧合。

把右端项改为 $x+y\ge b$，当 $b>0$ 时，同样推导得 $x^*(b)=y^*(b)=b/2$、$p^*(b)=b^2/4$、$\lambda^*(b)=b/2$。在 $b=1.2$ 处，右端项灵敏度为 $0.600000$，与中心差分一致。这里正乘子表示提高最低预算和会提高最优目标；若约束方向改变，符号解释也随之改变。

\MFQLead{非负约束真正改变解}

考虑
\[
 \Sigma=\begin{pmatrix}0.04&0.06\\0.06&0.16\end{pmatrix},
 \qquad \mathbf1^Tw=1.
\]
若只保留预算等式，公式
\[
 w_{\rm eq}=\frac{\Sigma^{-1}\mathbf1}{\mathbf1^T\Sigma^{-1}\mathbf1}
 =(1.25,-0.25)^T
\]
给出一个含空头的最小方差组合。加入 $w\ge0$ 后，这个点不可行；活跃集候选 $w_2=0$、$w_1=1$ 给出 $w_{\rm long}=(1,0)^T$，方差为 $0.04$。此例说明 long-only 不是一个无关紧要的报告选项：它改变可行域、活跃约束和最优性证书。数值求解后必须从候选权重重新检查这些条件，不能只比较目标值。

\section{数值证书与组合例子}

投影梯度迭代为
\[
x_{k+1}=\Pi_C\{x_k-\eta\nabla f(x_k)\}.
\]
对凸、梯度 $L$-Lipschitz 的目标，$0<\eta<2/L$ 是常见收敛区间；强凸还给线性收敛率。步长、尺度与停止规则不匹配会导致振荡或伪收敛。从 $(2,0)$ 出发、步长 $0.2$ 投影 500 次，本例数值解为 $(0.500000,0.500000)$。

核对不能只读 \texttt{success=True}，而应从候选解重新计算：最大原可行违反、负乘子违反、驻点范数、互补松弛和原始--对偶间隙。迭代差很小可能只是步长太小；目标变化很小也可能发生在平坦但不可行区域。

若 $f$ 的梯度是 $L$-Lipschitz，定义梯度映射
\[
 G_\eta(x)=\frac1\eta\{x-\Pi_C(x-\eta\nabla f(x))\}.
\]
投影的最优性与光滑性给出
\[
 f(x_{k+1})\le f(x_k)
 -\eta\left(1-\frac{L\eta}{2}\right)\|G_\eta(x_k)\|_2^2.
\]
因此 $0<\eta\le1/L$ 时目标单调下降；停止条件应控制 $\|G_\eta(x_k)\|$、可行违反和目标变化，而不是只看相邻迭代点。若 $f$ 还是 $m$-强凸，进一步可得到几何收敛；没有强凸性时，目标下降仍可能伴随变量在平坦方向漂移。

\MFQLead{量化例子：两资产最小方差组合}

给定
\[
\Sigma=\begin{pmatrix}0.04&0.006\\0.006&0.09\end{pmatrix},
\qquad \mathbf1^Tw=1,
\]
最小化 $w^T\Sigma w$。若暂不加非负约束，等式 KKT 给
\[
2\Sigma w-\nu\mathbf1=0,\qquad
w^*=\frac{\Sigma^{-1}\mathbf1}{\mathbf1^T\Sigma^{-1}\mathbf1}.
\]
本例最小方差权重为 $(0.711864,0.288136)$，方差为 $0.030203$，两项均为正，故也满足 long-only 约束。矩阵二范数条件数为 $2.308723$。把协方差非对角元从 $0.006$ 改为 $0.02$ 后，权重变为 $(0.777778,0.222222)$、方差为 $0.035556$、条件数为 $2.941260$；权重 $L^1$ 变化除以输入改变量为 $9.416196$。这是一项事先确定的扰动记录，而不是“权重看起来还行”的主观判断。

现实目标还可能加入线性/分段线性交易成本、二次冲击、换手约束、行业与风格暴露等。只要协方差半正定、成本凸且约束凸，问题仍可给全局证书；固定费用、最小手数和整数持仓会引入非凸或混合整数结构，不能沿用同一 KKT 全局结论。

\section{失败边界与研究审计}

非凸目标 $f(x)=(x^2-1)^2$ 在 $x=0$ 满足 $f'(x)=0$，但驻点 $x=0$ 的目标值为 1，而 $x=\pm1$ 的全局最优值为 0。一阶驻点不是全局证书。

问题 $\min_xx$、约束 $x^2\le0$ 的唯一可行点 $x=0$ 必然最优，但活跃约束梯度为零，驻点方程 $1+\lambda\cdot0=0$ 无解，KKT 驻点残差恒为 1。这里 Slater 与 LICQ 都失败，不能用“不存在乘子”否定最优点。

此外，$x\ge1$ 与 $x\le0$ 同时出现时不可行；无约束最小化 $-x$ 时目标无界。把同一约束乘以 $10^8$ 不改变精确可行域，却会改变数值残差尺度和停止判断。研究包必须区分数学状态（最优、不可行、无界）与数值状态（迭代上限、尺度不良、容差过宽）。

\MFQLead{研究审计协议}

提交优化结果时至少记录：变量与目标单位、凸性依据、数据估计窗口、可行域和活跃约束、求解器与容差、KKT/对偶证书、输入扰动敏感度、样本外成本及失败边界。对协方差、预期收益和成本系数做小扰动，比较权重、目标和活跃集；若微小输入导致仓位翻转，应报告不确定性或加入稳健化，而不是隐藏在更多小数位后。

\section{正则化为什么产生软阈值}

考虑一维凸目标 $f(x)=\tfrac12(x-a)^2+\lambda|x|$，$\lambda\geq0$。$x>0$ 时驻点为 $x=a-\lambda$，只有 $a>\lambda$ 时落在该区域；$x<0$ 时驻点为 $x=a+\lambda$，需要 $a<-\lambda$。在零点，$0\in-a+\lambda[-1,1]$ 等价于 $|a|\leq\lambda$。合并得到
\[
 x^*=\operatorname{sign}(a)(|a|-\lambda)_+.
\]
平方项强凸，故这个解唯一。$a=0.8,\lambda=0.3$ 时解为 $0.5$；当 $a$ 降至 $0.2$ 时，解直接为零。notebook 将逐段计算这个解，并与目标曲线的最低点比较。稀疏性来自零点处一整段次梯度，不能通过把绝对值当成处处可微来理解。

\noindent\textbf{延伸阅读。}凸分析、对偶与 KKT 参见\cite{boyd2004convex}，数值优化实现边界参见\cite{nocedal2006numerical}。

\section{四级问题与即时练习}
\begingroup\small
\begin{enumerate}[nosep,leftmargin=*]
  \item \textbf{口述概念：}区分凸集、凸函数、约束资格、KKT 必要性、充分性与强对偶。
  \item \textbf{笔试推导：}手算二次规划的 KKT 解、对偶函数与零对偶间隙。
  \item \textbf{数值编程：}实现投影梯度并独立报告四类 KKT 残差与对偶间隙，而非只读优化器状态。
  \item \textbf{研究判断：}带换手、成本与暴露约束的组合优化结果极不稳定，应先审计哪些建模与数值边界？
\end{enumerate}
\endgroup

\subsection*{分级提示}
\begingroup\small
\begin{enumerate}[nosep,leftmargin=*]
  \item \textbf{口述概念：}先区分问题是否凸，再说明 KKT 作为必要或充分条件还缺什么假设。
  \item \textbf{笔试推导：}先判断哪条约束活跃；从驻点得到 $x=y=\lambda$，再代入互补松弛。
  \item \textbf{数值编程：}投影保持可行；验收值来自手算解，残差由候选解重新计算。
  \item \textbf{研究判断：}核对单位、凸性、可行性、条件数、活跃集、参数扰动与样本外成本。
\end{enumerate}
\endgroup


\noindent\textbf{本章要点。}凸性决定局部证据能否成为全局证书，约束资格决定乘子是否存在，对偶间隙与灵敏度提供可计算审计。一个可复核的优化结果至少要同时报告单位、可行性、活跃集、驻点残差、对偶间隙、输入扰动和失败案例。
